\subsection{\texorpdfstring{\(D_{l}\) 型系统 (\(l \geqslant 3\))}{D\_l 型系统 (l >= 3)}}

\begin{enumerate}
\def\labelenumi{(\Roman{enumi})}
\item
  在 \(V = R^{l}\) 中考虑第 4 号的群 \(L_{0}\)。集合 R 中满足 \(\alpha \in L_{0}\) 且 \((\alpha|\alpha) = 2\) 的根由向量 \(\pm\varepsilon_{i} \pm \varepsilon_{j}\)（ \(1 \leqslant i < j \leqslant l\)）组成。显然，R 生成 V，且 \(2(\alpha|\beta)/(\alpha|\alpha) \in \mathbf{Z}\) 对所有 \(\alpha, \beta \in R\) 成立。因此，R 是 V 中的约化根系（第 4 号）。根的数目为 \(n = 2l(l - 1)\)。
\item
  令
\end{enumerate}

\[
\alpha_ {1} = \varepsilon_ {1} - \varepsilon_ {2}, \alpha_ {2} = \varepsilon_ {2} - \varepsilon_ {3} \dots . \alpha_ {l - 1} = \varepsilon_ {l - 1} - \varepsilon_ {l}, \alpha_ {l} = \varepsilon_ {l - 1} + \varepsilon_ {l}.
\]

下列公式是显然的：

\[
\begin{array}{c} \varepsilon_ {i} - \varepsilon_ {j} = \alpha_ {i} + \alpha_ {i + 1} + \dots + \alpha_ {j - 1} (i <   j) \\ \varepsilon_ {i} + \varepsilon_ {j} = \alpha_ {i} + \alpha_ {i + 1} + \dots + \alpha_ {j - 1} + 2 \alpha_ {j} + 2 \alpha_ {j + 1} + \dots \\ \qquad \qquad \qquad \qquad \qquad \qquad \dots + 2 \alpha_ {l - 2} + \alpha_ {l - 1} + \alpha_ {l} (i <   j \leqslant l - 2) \\ \varepsilon_ {i} + \varepsilon_ {l - 1} = \alpha_ {i} + \alpha_ {i + 1} + \dots + \alpha_ {l} (i <   l - 1) \\ \varepsilon_ {i} + \varepsilon_ {l} = \alpha_ {i} + \alpha_ {i + 1} + \dots + \alpha_ {l - 2} + \alpha_ {l} (i <   l - 1) \\ \varepsilon_ {l - 1} + \varepsilon_ {l} = \alpha_ {l}. \end{array}
\]

所以 \((\alpha_{1},\ldots,\alpha_{l})\) 是 R 的一个基（§ 1，第 2 号，命题 20 的推论 3）。此外，\(\| \alpha_{i}\|^{2}=2\) 对所有 \(i\) 成立，且 \((\alpha_{i}|\alpha_{j})=0\) 对 \(i+1<j\) 成立，但 \(i=l-2\)、\(j=l\) 时例外，此时 \((\alpha_{l-2}|\alpha_{l})=-1\)；\((\alpha_{i}|\alpha_{i+1})=-1\) 对 \(i\leqslant l-2\) 成立，最后 \((\alpha_{l-1}|\alpha_{l})=-1\)；因此 R 的 Dynkin 图为 D \(_{l}\) 型。正根是 \(\varepsilon_{i}\pm\varepsilon_{j}\)，其中 \(i<j\)。

\begin{enumerate}
\def\labelenumi{(\Roman{enumi})}
\setcounter{enumi}{2}
\item
  我们有 \(h = n / l = 2(l - 1)\)。
\item
  令 \(\tilde{\alpha} = \varepsilon_{1} + \varepsilon_{2} = \alpha_{1} + 2\alpha_{2} + \cdots + 2\alpha_{l-2} + \alpha_{l-1} + \alpha_{l}\)，这是一个根。相对于 \((\alpha_{i})\)，它的坐标之和为
\end{enumerate}

\[
2 l - 3 = h - 1.
\]

因此 \(\tilde{\alpha}\) 是最高根。

若 \(l = 3\)，则有

\[
(\tilde {\alpha} | \alpha_ {1}) = 0, \quad (\tilde {\alpha} | \alpha_ {2}) = (\tilde {\alpha} | \alpha_ {3}) = 1.
\]

若 \(l \geqslant 4\)，则 \((\tilde{\alpha}|\alpha_i) = 0\) 当 \(i \neq 2\) 时成立，且 \((\tilde{\alpha}|\alpha_2) = 1\)。因此，扩展 Dynkin 图为：

\[
\alpha_ {1} \quad \alpha_ {2} \quad \alpha_ {3} \quad \dots \quad \alpha_ {i - 3} \quad \alpha_ {i - 2} \quad \alpha_ {i - 1}
\]

\begin{enumerate}
\def\labelenumi{(\Alph{enumi})}
\setcounter{enumi}{21}
\tightlist
\item
  由于 \((\alpha |\alpha) = 2\) 对所有 \(\alpha \notin R\) 成立，所以 \(R^{\sim} = R\)。
\end{enumerate}

对于 \(\Phi_{\mathrm{R}}\)，根的长度为 \(h^{-1 / 2} = (2l - 2)^{-1 / 2}\)。因此

\[
\Phi_ {R} (x, y) = (x | y) / (4 l - 4) \quad \text { and } \quad \gamma (R) = 4 (l - 1) ^ {2}.
\]

\begin{enumerate}
\def\labelenumi{(\Roman{enumi})}
\setcounter{enumi}{5}
\tightlist
\item
  类似于第 7 号的计算给出基本权：
\end{enumerate}

\[
\begin{array}{r l} \omega_ {i} & = \varepsilon_ {1} + \varepsilon_ {2} + \dots + \varepsilon_ {i} \\ & = \alpha_ {1} + 2 \alpha_ {2} + \dots + (i - 1) \alpha_ {i - 1} + i (\alpha_ {i} + \alpha_ {i + 1} + \dots + \alpha_ {l - 2}) \\ & \quad + \frac {1}{2} i (\alpha_ {l - 1} + \alpha_ {l}) \end{array}
\]

当 i \textless{} l - 1 时，

\[
\begin{array}{l} \omega_ {l - 1} = \frac {1}{2} (\varepsilon_ {1} + \varepsilon_ {2} + \dots + \varepsilon_ {l - 2} + \varepsilon_ {l - 1} - \varepsilon_ {l}) \\ \qquad = \frac {1}{2} (\alpha_ {1} + 2 \alpha_ {2} + \dots + (l - 2) \alpha_ {l - 2} + \frac {1}{2} l \alpha_ {l - 1} + \frac {1}{2} (l - 2) \alpha_ {l}), \\ \omega_ {l} = \frac {1}{2} (\varepsilon_ {1} + \varepsilon_ {2} + \dots + \varepsilon_ {l - 2} + \varepsilon_ {l - 1} + \varepsilon_ {l}) \\ \qquad = \frac {1}{2} (\alpha_ {1} + 2 \alpha_ {2} + \dots + (l - 2) \alpha_ {l - 2} + \frac {1}{2} (l - 2) \alpha_ {l - 1} + \frac {1}{2} l \alpha_ {l}). \end{array}
\]

\begin{enumerate}
\def\labelenumi{(\Roman{enumi})}
\setcounter{enumi}{6}
\tightlist
\item
  正根之和为
\end{enumerate}

\[
\begin{array}{l} 2 \rho = 2 (l - 1) \varepsilon_ {1} + 2 (l - 2) \varepsilon_ {2} + \dots + 2 \varepsilon_ {l - 1} \\ = \sum_ {i = 1} ^ {l - 2} 2 (i l - \frac {i (i + 1)}{2}) \alpha_ {i} + \frac {l (l - 1)}{2} (\alpha_ {l - 1} + \alpha_ {l}). \end{array}
\]

\begin{enumerate}
\def\labelenumi{(\Roman{enumi})}
\setcounter{enumi}{7}
\item
  \(\pm \varepsilon_{i} \pm \varepsilon_{j}\) 生成 \(L_{1}\)（第 4 号），所以 \(Q(R) = L_{1}\)。因此 \(Q(R^{-}) = L_{1}\)，从而 \(P(R) = L_{2}\)（第 4 号）。由第 4 号，\(P(R)/Q(R)\) 同构于 \(\mathbf{Z}/4\mathbf{Z}\)（当 \(l\) 为奇数时），同构于 \(\mathbf{Z}/2\mathbf{Z} \times \mathbf{Z}/2\mathbf{Z}\)（当 \(l\) 为偶数时）。在第一种情形中，\(P(R)/Q(R)\) 由 \(\omega_{l}\) 的规范像生成（也由 \(\omega_{l-1}\) 的规范像生成）；在第二种情形中，\(P(R)/Q(R)\) 由 \(\omega_{l-1}\) 和 \(\omega_{l}\) 的规范像生成。在两种情形中，连接指标都是 4。
\item
  以及 (X) 在 \(\mathbf{R}^l\) 中，正交反射 \(s_{\varepsilon_i - \varepsilon_j}\)（ \(i \neq j\)）将 \(\varepsilon_i\) 与 \(\varepsilon_j\) 互换，并保持 \(\varepsilon_k\) 不变，其中 \(k\) 与 \(i\) 、 \(j\) 均不同。\(s_{\varepsilon_i - \varepsilon_j}\) 生成群 \(\mathbf{G}_1\)，它同构于对称群 \(\mathfrak{S}_l\)。另一方面，\(s_{ij} = s_{\varepsilon_i - \varepsilon_j} s_{\varepsilon_i + \varepsilon_j}\) 将 \(\varepsilon_i\) 变为 \(-\varepsilon_i\)，将 \(\varepsilon_j\) 变为 \(-\varepsilon_j\)，并保持 \(\varepsilon_k\) 不变，其中 \(k\) 与 \(i\) 、 \(j\) 均不同。\(s_{ij}\) 生成群 \(\mathbf{G}_2\)，它是由自同构 \(u\) 组成的 \(\mathbf{R}^l\) 的自同构群，满足 \(u(\varepsilon_i) = (-1)^{\nu_i} \varepsilon_i\) 且 \(\prod_{i=1}^{l} (-1)^{\nu_i} = 1\)。群 \(\mathbf{G}_2\) 同构于 \((\mathbf{Z}/2\mathbf{Z})^{l-1}\)，且 \(\mathbf{G}_2\) 是 W(R) 的正规子群，所以 W(R) 同构于 \(\mathfrak{S}_l\) 与 \((\mathbf{Z}/2\mathbf{Z})^{l-1}\) 的半直积。因此，其阶为 \(2^{l-1}.l!\)。
\end{enumerate}

 第 5 号的多项式函数 \(t_j\) 在 \(W(R)\) 下不变，\(t = \xi_1\xi_2\ldots\xi_l\) 也不变；此外，\(t_l = t^2\)。设 \(P(\xi_1,\ldots ,\xi_l)\) 是在 \(W(R)\) 下不变的多项式。设单项式 \(\xi_1^{\nu_1}\xi_2^{\nu_2}\ldots\xi_l^{\nu_l}\) 出现在 \(P\) 中且满足 \(\nu_{i}\) 为奇数；那么 \(\nu_{j}\) 对所有 \(j\) 都是奇数，因为单项式 \((-1)^{\nu_i + \nu_j}\xi_1^{\nu_1}\xi_2^{\nu_2}\ldots\xi_l^{\nu_l}\) 出现在 \(s_{ij}(P)\) 中，所以 \(\nu_{i} + \nu_{j} \equiv 0 \pmod{2}\)，且 \(\nu_{j} \equiv 1 \pmod{2}\)。于是

\(P = P_{1} + tP_{2}\)，其中出现在 \(P_{1}\) 和 \(P_{2}\) 中的所有单项式都只有偶数指数。由于 P 在 \(\xi_{i}\) 的置换下不变，\(P_{1}\) 和 \(P_{2}\) 也具有同样的性质，因此可以写成 \(t_{1}, t_{2}, \ldots, t_{l}\) 的多项式。这证明代数 \(\mathbf{S}(\mathbf{R}^{l})^{\mathrm{W}(\mathbb{R})}\) 由 \(t_{1}, t_{2}, \ldots, t_{l-1}, t\) 生成。此外，\(\mathbf{S}(\mathbf{R}^{l})^{\mathrm{W}(\mathbb{R})}\) 的分式域的超越次数为 l，所以 \(t_{1}, t_{2}, \ldots, t_{l-1}, t\) 代数独立。我们得到，适当排序后的指数序列为：

\[
1, 3, 5, \dots , 2 l - 5, 2 l - 3, l - 1.
\]

注意，\(l - 1\) 在 \(l\) 为偶数时出现两次，在 \(l\) 为奇数时出现一次。

\begin{enumerate}
\def\labelenumi{(\Roman{enumi})}
\setcounter{enumi}{10}
\tightlist
\item
  Dynkin 图的自同构就是其底层图的自同构。因此：
\end{enumerate}

\begin{enumerate}
\def\labelenumi{\arabic{enumi})}
\item
  若 \(l = 3\)，则 A(R)/W(R) 同构于 Z/2Z。
\item
  若 \(l = 4\)，则终端顶点的每个置换都定义图的一个自同构，所以 \(\mathrm{A}(\mathrm{R}) / \mathrm{W}(\mathrm{R})\) 同构于 \(\mathfrak{S}_3\)。
\item
  若 \(l \geqslant 5\)，则从分叉点出发的链长度为 1、1 和 \(l - 3 \geqslant 2\)。因此，除恒等自同构外，图的唯一自同构对应于自同构 \(\varepsilon \in \mathrm{A}(\mathrm{R})\)：它互换 \(\alpha_{l-1}\) 与 \(\alpha_l\)，并固定 \(\alpha_i\)（\(1 \leqslant i \leqslant l - 2\)）。于是 \(\mathrm{A}(\mathrm{R}) / \mathrm{W}(\mathrm{R})\) 同构于 \(\mathbf{Z}/2\mathbf{Z}\)；此外，\(\mathrm{A}(\mathrm{R})\) 是 (IX) 中定义的群 \(\mathbf{G}_1 \cong \mathfrak{S}_l\) 与群 \(\mathbf{G}_3\) 的半直积；该群由自同构 \(u\) 组成，这些自同构是 \(\mathbf{R}^l\) 的自同构，且 \(u(\varepsilon_i) = \pm \varepsilon_i\) 对所有 \(i\) 成立。
\end{enumerate}

若 \(l\) 为偶数，则 \(-1 \in \mathrm{W}(\mathrm{R})\)，所以 \(w_0 = -1\)。若 \(l\) 为奇数，则 \(-1 \notin \mathrm{W}(\mathrm{R})\)，所以 \(\mathrm{A}(\mathrm{R}) = \mathrm{W}(\mathrm{R}) \times \{1, -1\}\)，且 \(w_0 = -\varepsilon\)。

\begin{enumerate}
\def\labelenumi{(\Roman{enumi})}
\setcounter{enumi}{11}
\tightlist
\item
当 \(l\) 为偶数时，\(\mathrm{P}(\mathbb{R}^{-}) / \mathrm{Q}(\mathbb{R}^{-})\) 有三个 2 阶元素，即 \(\omega_{1}, \omega_{l-1}\) 和 \(\omega_{l}\)。由于 \(\omega_{l}\)（分别为 \(\omega_{l-1}\)）互换对应于 \(\alpha_{0}\) 和 \(\alpha_{l}\)（分别为 \(\alpha_{l-1}\)）的顶点，它也互换对应于 \(\alpha_{1}\) 和 \(\alpha_{l-1}\)（分别为 \(\alpha_{l}\)）的顶点，并且还互换对应于 \(\alpha_{j}\) 和 \(\alpha_{l-j}\) 的顶点，其中
\end{enumerate}

\[
2 \leqslant j \leqslant l - 2.
\]

我们有 \(\omega_{1} = \omega_{l}\omega_{l - 1}\)

当 \(l\) 为奇数时，\(\mathrm{P}(\mathbb{R}^{-}) / \mathrm{Q}(\mathbb{R}^{-})\) 有两个 4 阶元素，即 \(\omega_{l-1}\) 和 \(\omega_l\)，以及一个 2 阶元素，即 \(\omega_1\)。事实上，\(\omega_1\) 互换对应于 \(\alpha_0\) 和 \(\alpha_1\) 的顶点，所以它固定对应于 \(\alpha_j\) 的顶点，其中 \(2 \leqslant j \leqslant l-2\)，并且必为 2 阶。因此，\(\omega_l\) 是 4 阶的，并将对应于 \(\alpha_0\)（分别为 \(\alpha_l\)、\(\alpha_1\)、\(\alpha_{l-1}\)）的顶点变为对应于 \(\alpha_l\)（分别为 \(\alpha_1\)、\(\alpha_{l-1}\)、\(\alpha_0\)）的顶点，并互换对应于 \(\alpha_j\) 和 \(\alpha_{l-j}\) 的顶点，其中 \(2 \leqslant j \leqslant l-2\)。我们有 \(\omega_1 = \omega_l^2\)，且 \(\omega_{l-1} = \omega_l^3\)。

当 \(l \neq 4\) 时，\(\mathrm{A}(\mathrm{R})/\mathrm{W}(\mathrm{R})\) 的非单位元互换对应于 \(\alpha_{l-1}\) 和 \(\alpha_{l}\) 的顶点，因而互换元素 \(\omega_{l-1}\) 和 \(\omega_{l}\)（它们属于 \(\mathrm{P}(\mathrm{R}^{-})/\mathrm{Q}(\mathrm{R}^{-})\)）。当 l 为奇数时，由此得到的 \(\mathrm{P}(\mathrm{R}^{-})/\mathrm{Q}(\mathrm{R}^{-})\) 的自同构是映射 \(x \mapsto -x\)。

当 \(l = 4\) 时，\(A(R)/W(R)\) 可识别为 \(\{1,3,4\}\) 的置换群，并通过对指标的置换作用于 \(\{\omega_1,\omega_3,\omega_4\}\)。

